% ===================================================================
\section{History: people, experiments, ideas}\label{sec:history}
% ===================================================================
\begin{trapbox}\textbf{Verification status:} every entry below is editor knowledge.\src{EK} No primary or scholarly source could be accessed in this build (network-restricted), so dates are given at the precision the editor is confident of; ``c.'' marks approximate dates. Confirm any date before relying on it to break a tie. Types: \textbf{E} experiment/observation, \textbf{P} publication/theory, \textbf{I} invention/patent, \textbf{N} later naming.\end{trapbox}

\begin{reftab}{T:history}{Compact timeline (chronological).}
\small
\renewcommand{\arraystretch}{1.08}
\begin{tabularx}{\textwidth}{@{}lL{0.21\textwidth}cY@{}}
\toprule
Year & Person(s) & Type & Contribution and connection\\
\midrule
c.\,1593 & Galileo Galilei (attribution debated) & I & Air thermoscope: shows hotter/colder without a scale; precursor of the thermometer.\\
1714; 1724 & D.\,G. Fahrenheit & I, P & Mercury-in-glass thermometer; published his temperature scale (1724).\\
1742 & Anders Celsius & P & Centigrade scale with \emph{100 at the freezing point and 0 at boiling}; inverted to the modern orientation shortly after his death.\\
c.\,1760s & Joseph Black & E & Latent heat and specific heat (``heat capacity''); distinguished quantity of heat from temperature. Lectures published posthumously (1803).\\
1780s & A.\,Lavoisier \& P.-S. Laplace & E & Ice calorimeter: measured heat by mass of ice melted. Lavoisier's caloric (heat-fluid) theory.\\
1798 & Count Rumford (B. Thompson) & E & Cannon boring produced seemingly inexhaustible heat from friction; challenged caloric theory.\\
1822 & Joseph Fourier & P & \emph{Analytical Theory of Heat}: conduction law $q''=-k\nabla T$ and the heat equation.\\
1824 & Sadi Carnot & P & \emph{Reflections on the Motive Power of Fire}: ideal reversible engine; efficiency limited by reservoir temperatures.\\
1834 & \'Emile Clapeyron & P & Put Carnot's ideas in analytic form with the $P$--$V$ diagram; Clapeyron relation (later completed by Clausius).\\
1842--1850 & J.\,R. Mayer; J.\,P. Joule; H. von Helmholtz & E, P & Mechanical equivalent of heat (Joule's paddle-wheel experiments, 1840s); conservation of energy (Mayer 1842, Helmholtz 1847).\\
1848 & William Thomson (Lord Kelvin) & P & Absolute temperature scale independent of any substance.\\
1850; 1865 & Rudolf Clausius & P, N & Clear statements of the first and second laws (1850); coined ``entropy'' and wrote $\dd S=\delta Q_{\text{rev}}/T$ (1865).\\
1852 & Joule \& Thomson & E & Porous-plug (throttling) experiments: Joule--Thomson effect.\\
1859 & W.\,J.\,M. Rankine & P & \emph{Manual of the Steam Engine}; the vapor power cycle now named for him.\\
1860; 1872--77 & J.\,C. Maxwell; L. Boltzmann & P & Velocity distribution (Maxwell 1860); statistical entropy and the $H$-theorem (Boltzmann). The form $S=k\ln W$ and the symbol $k$ were written by Planck (c.\,1900).\\
1873 & J.\,D. van der Waals & P & Doctoral thesis: equation of state with attraction and excluded volume; continuity of liquid and gas.\\
1875--1878 & J.\,W. Gibbs & P & \emph{On the Equilibrium of Heterogeneous Substances}: chemical potential, phase rule, free energy.\\
1876; 1890s & N. Otto; R. Diesel & I & Four-stroke spark-ignition engine (Otto); compression-ignition engine (Diesel patent, 1890s).\\
1879; 1884 & J. Stefan; L. Boltzmann & E, P & Stefan's empirical $T^4$ law of radiation; Boltzmann's thermodynamic derivation.\\
1893; 1900 & W. Wien; M. Planck & P & Displacement law; Planck's radiation law with energy quanta.\\
1906; 1911 & W. Nernst; Planck & P & Nernst heat theorem $\to$ third law; Planck's $S\to0$ form.\\
1908 & H. Kamerlingh Onnes & E & First liquefaction of helium (using precooling plus Joule--Thomson expansion).\\
2019 & CGPM (SI redefinition) & N & Kelvin defined by fixing $k=\SI{1.380649e-23}{J/K}$ exactly.\src{CODATA}\\
\bottomrule
\end{tabularx}\end{reftab}
\textbf{Shared attribution:} energy conservation emerged from several workers (Mayer, Joule, Helmholtz, Colding); ``Boyle's law'' was also found by Mariotte; ``Charles's law'' was published by Gay-Lussac. Credit questions should be answered with the names the question's source uses.\src{EK}
